% !TEX TS-program = xelatex
\documentclass[10pt,a4paper]{article}

\usepackage[a4paper, left=20mm, right=20mm, top=15mm, bottom=15mm]{geometry}
\usepackage{fontspec}
\setmainfont{TeX Gyre Heros}
\usepackage{amsmath,amssymb,mathtools}
\usepackage{array,booktabs}
\usepackage{enumitem}
\usepackage{mdframed}
\usepackage{xcolor}
\usepackage{tikz}
\usetikzlibrary{calc}
\usepackage{eso-pic}
\usepackage[useregional]{datetime2}

% Watermark
\newcommand{\mymarginlabel}{\textcopyright\ 2026 David Ewing dew59@uclive.ac.nz | \DTMnow\ | \jobname.tex}
\AddToShipoutPictureBG{%
  \begin{tikzpicture}[remember picture,overlay]
    \node[rotate=90, anchor=north]
      at ($(current page.north west)+(1.4cm,-13cm)$)
      {\scriptsize\ttfamily \mymarginlabel};
  \end{tikzpicture}%
}

\definecolor{reasonbg}{RGB}{240,248,240}
\definecolor{boxbg}{RGB}{235,242,252}

\setlength{\parskip}{3pt}
\setlength{\parindent}{0pt}
\setlength{\abovedisplayskip}{5pt}
\setlength{\belowdisplayskip}{5pt}
\setlength{\abovedisplayshortskip}{3pt}
\setlength{\belowdisplayshortskip}{3pt}

% Section heading with rule and mark count
\newcommand{\sectitle}[2]{%
  \vspace{10pt}\noindent\rule{\linewidth}{0.8pt}\par\vspace{2pt}%
  {\large\textbf{#1}\hfill\normalsize\textit{[#2 marks]}}\par
  \vspace{2pt}\rule{\linewidth}{0.4pt}\vspace{4pt}}

% Green reasoning box
\newcommand{\reason}[1]{%
  \begin{mdframed}[backgroundcolor=reasonbg, linewidth=0.4pt,
    innertopmargin=3pt, innerbottommargin=3pt,
    innerleftmargin=6pt, innerrightmargin=6pt]
    \small\textit{Reason:} \small #1
  \end{mdframed}\vspace{2pt}}

% Blue result box
\newenvironment{resultbox}{%
  \begin{mdframed}[backgroundcolor=boxbg, linewidth=0.6pt,
    innertopmargin=5pt, innerbottommargin=5pt,
    innerleftmargin=8pt, innerrightmargin=8pt]
}{\end{mdframed}}

\begin{document}

\begin{center}
  {\large\textbf{ENEL445 --- Worked Examples: Quadratic Optimisation \& Optimality Conditions}}\\[3pt]
  {\small David Ewing (82171165)\quad|\quad
   Covers: variable transformation, optimality condition, Van~der~Corput,
   quadratic gradient, exact line search, optimal step size.}
\end{center}
\noindent\rule{\linewidth}{0.8pt}

% ==============================================================
\sectitle{Q1. Design Variable Transformation via Sigmoid}{5}
% ==============================================================

\noindent\textbf{Problem.} The $i$-th design variable $x_i$ must satisfy $|x_i|\le a$.
Eliminate this constraint using the sigmoid $\sigma(t)=\dfrac{1}{1+e^{-t}}$.

\vspace{6pt}
\textbf{Step 1 --- restate the constraint.}

$|x_i|\le a$ is equivalent to $x_i\in[-a,\,a]$, an interval of width $2a$.
\[
  -a \;<\; x_i \;<\; +a
\]

\textbf{Step 2 --- range of the sigmoid.}
\[
  \sigma(t)=\frac{1}{1+e^{-t}},\qquad t\in\mathbb{R}\;\Rightarrow\;\sigma(t)\in(0,1).
\]

\reason{If $t\to+\infty$, $e^{-t}\to 0$ so $\sigma\to 1$.
        If $t\to-\infty$, $e^{-t}\to\infty$ so $\sigma\to 0$.
        The image is strictly $(0,1)$ for all finite $t$.}

\textbf{Step 3 --- scale and shift to $(-a,a)$.}

Multiply by $2a$ to obtain range $(0,2a)$, then shift by $-a$:
\begin{align*}
  2a\,\sigma(t_i)             &\in (0,\,2a) \\
  x_i = -a + 2a\,\sigma(t_i) &\in (-a,\,+a).
\end{align*}

\reason{Sanity check: $t_i\to-\infty \Rightarrow x_i = -a+0 = -a$.
        $t_i\to+\infty \Rightarrow x_i = -a+2a = +a$.
        For all finite $t_i$, $x_i\in(-a,+a)\subseteq[-a,a]$. \checkmark}

\textbf{Step 4 --- explicit bound confirmation.}

For all finite $t_i\in\mathbb{R}$, the transformed variable satisfies:
\[
  x_i \in (-a,\,+a).
\]
The open interval $(-a,+a)$ is \textbf{always} a subset of the closed constraint set:
\[
  (-a,+a) \subseteq [-a,\,a]. \checkmark
\]
\reason{The sigmoid never reaches its limits for finite $t_i$, so the image is strictly
        the open interval $(-a,+a)$, giving margin on both sides. The constraint is
        therefore satisfied automatically and can be dropped entirely from the optimisation.}

\newpage
\textbf{Why use $\sigma(t)$?}

The sigmoid is one member of a family of functions that share three key properties
required for this variable-transformation trick:
\begin{enumerate}
  \item \textbf{Bounded range} --- maps $\mathbb{R}$ into a finite interval,
        so the constraint on $x_i$ is encoded implicitly.
  \item \textbf{Smooth and differentiable everywhere} --- so the chain rule
        $\partial f/\partial t_i = (\partial f/\partial x_i)(\partial x_i/\partial t_i)$
        is well defined, and gradient-based optimisers can be used on $t_i$
        without modification.
  \item \textbf{Strictly monotone} --- the mapping $t_i\mapsto x_i$ is one-to-one,
        so no information is lost and every $x_i\in(-a,a)$ has a unique pre-image.
\end{enumerate}

\reason{Any smooth, strictly monotone function with bounded range works.
        Alternatives include $\tanh(t)$ (range $(-1,1)$, scaled by $a$) and
        $\tfrac{2}{\ \pi}\arctan(t)$ (range $(-1,1)$).
        The sigmoid is convenient because its range is exactly $(0,1)$,
        which maps cleanly to $(-a,a)$ with one multiply-and-shift.
        Without this trick, a constrained method (e.g.\ KKT conditions)
        would be needed; the transformation converts the problem into
        an unconstrained one solvable by any standard solver.}

\begin{resultbox}
Substitute $\displaystyle x_i = -a + \frac{2a}{1+e^{-t_i}}$
and optimise over unconstrained $t_i\in\mathbb{R}$.
The constraint $|x_i|\le a$ is then automatically satisfied and can be dropped.
\smallskip\\
\textit{Equivalent forms:}\;
$x_i = a\!\left(\dfrac{2}{1+e^{-t_i}}-1\right) = a\tanh\!\left(\dfrac{t_i}{2}\right)$.
\end{resultbox}

\newpage
% ==============================================================
\sectitle{Q2. First-Order Necessary Optimality Condition}{5}
% ==============================================================

\noindent\textbf{Problem.} Show that if $x^*$ is a local minimum of $f(x)$,
then $\nabla f(x^*)=0$.
\textit{Hint: use the first-order Taylor expansion of $f(x)$.}

\vspace{6pt}
\textbf{Step 1 --- local-minimum definition.}
There exists $\delta>0$ such that
\[
  f(x^*+d)\ge f(x^*)\quad\text{for all }\|d\|<\delta.
\]

\textbf{Step 2 --- first-order Taylor expansion about $x^*$.}
\begin{align*}
  f(x)     &= f(x^*) + \nabla f(x^*)^\top(x-x^*) + O(\|x-x^*\|^2) \\
  f(x^*+d) &= f(x^*) + \nabla f(x^*)^\top d       + O(\|d\|^2).
\end{align*}

\reason{The chain rule gives the linear term $\nabla f(x^*)^\top d$;
        higher-order terms are $O(\|d\|^2)$.}

\textbf{Step 3 --- substitute into the local-minimum inequality.}

Substituting the Taylor expansion into $f(x^*+d)\ge f(x^*)$:
\begin{align*}
  f(x^*) + \nabla f(x^*)^\top d + O(\|d\|^2) &\ge f(x^*) \\
  \nabla f(x^*)^\top d + O(\|d\|^2)           &\ge 0\quad\text{for all }\|d\|<\delta.
\end{align*}

\textbf{Step 4 --- choose a test direction (contradiction).}
Set $d=-\varepsilon\,\nabla f(x^*)$ with $\varepsilon>0$ small:
\[
  -\varepsilon\,\|\nabla f(x^*)\|^2 + O(\varepsilon^2)\ge 0.
\]

\textbf{Step 5 --- divide by $\varepsilon$ and let $\varepsilon\to 0^+$.}
\reason{Dividing by $\varepsilon>0$ preserves the inequality; the $O(\varepsilon)$ term vanishes.}
\begin{align*}
  -\|\nabla f(x^*)\|^2 &\ge 0 \\
   \|\nabla f(x^*)\|^2 &\le 0
\intertext{But a squared norm is \textbf{always} non-negative:}
   \|\nabla f(x^*)\|^2 &\ge 0
\end{align*}
\begin{resultbox}
\[
  0 \;\le\; \|\nabla f(x^*)\|^2 \;\le\; 0
  \;\implies\; \|\nabla f(x^*)\|^2 = 0
  \;\implies\; \nabla f(x^*)=0. \hfill\square
\]
\end{resultbox}

%\newpage
% ==============================================================
\sectitle{Q3. Van der Corput Sequence, base $b=3$}{10}
% ==============================================================

\noindent\textbf{Definition.} If $i=a_0 + a_1 b + a_2 b^2+\cdots+a_r b^r$ in base $b$, then
\[
  \phi_i(b) = \frac{a_0}{b}+\frac{a_1}{b^2}+\cdots+\frac{a_r}{b^{r+1}}.
\]

\textbf{Algorithm (three steps for each $i$):}
\begin{enumerate}[noitemsep]
  \item Write $i$ in base $b$, reading off digits $a_0, a_1, \ldots, a_r$ least-significant first.
  \item Each digit $a_k$ contributes $a_k/b^{k+1}$ to the sum.
  \item Sum all contributions: $\phi_i(b) = \displaystyle\sum_{k=0}^{r} \dfrac{a_k}{b^{k+1}}$.
\end{enumerate}

\vspace{4pt}
\begin{center}
\begin{tabular}{cclll}
\toprule
$i$ & Base-3 repr. & Digits $a_0, a_1, \ldots$ & Computation & $\phi_i(3)$ \\
\midrule
1 & $(1)_3$  & $a_0=1$           & $1/3$                 & $1/3$ \\
2 & $(2)_3$  & $a_0=2$           & $2/3$                 & $2/3$ \\
3 & $(10)_3$ & $a_0=0,\;a_1=1$   & $0/3+1/9$             & $1/9$ \\
4 & $(11)_3$ & $a_0=1,\;a_1=1$   & $1/3+1/9$             & $4/9$ \\
5 & $(12)_3$ & $a_0=2,\;a_1=1$   & $2/3+1/9$             & $7/9$ \\
\bottomrule
\end{tabular}
\end{center}

\reason{%
  $i=1$: $(1)_3$, $a_0=1 \Rightarrow 1/3$. \quad
  $i=2$: $(2)_3$, $a_0=2 \Rightarrow 2/3$.\\[2pt]
  $i=3$: $3=(1{\cdot}3+0)=(10)_3$, $a_0=0,\,a_1=1 \Rightarrow 0/3+1/9=1/9$.\\[2pt]
  $i=4$: $4=(1{\cdot}3+1)=(11)_3$, $a_0=1,\,a_1=1 \Rightarrow 1/3+1/9=4/9$.\\[2pt]
  $i=5$: $5=(1{\cdot}3+2)=(12)_3$, $a_0=2,\,a_1=1 \Rightarrow 2/3+1/9=7/9$.}

\begin{resultbox}
Van der Corput sequence in base $b=3$ for $i=1,\ldots,5$:
\begin{align*}
  \phi_1(3) &= \tfrac{1}{3} \\
  \phi_2(3) &= \tfrac{2}{3} \\
  \phi_3(3) &= \tfrac{1}{9} \\
  \phi_4(3) &= \tfrac{4}{9} \\
  \phi_5(3) &= \tfrac{7}{9}.
\end{align*}
\end{resultbox}

\newpage
% ==============================================================
\sectitle{Q4. Gradient of Quadratic Objective}{3}
% ==============================================================

\noindent\textbf{Problem.} Let $f(x)=\tfrac{1}{2}x^\top Ax + b^\top x$,
$A$ symmetric positive definite. Find $\nabla f(c)$.

\textbf{Linear term:} $\partial(b^\top x)/\partial x = b$.

\reason{Standard result: $\partial(a^\top x)/\partial x = a$.}

\textbf{Quadratic term} (using $A=A^\top$):
\[
  \frac{\partial}{\partial x}\!\left(\tfrac{1}{2}x^\top Ax\right)
  = \tfrac{1}{2}(A+A^\top)x = Ax.
\]

\reason{Product rule gives $(A+A^\top)x/2$; symmetry collapses this to $Ax$.}

\begin{resultbox}
$\nabla f(x) = Ax + b$,\quad so\quad $\nabla f(c) = Ac + b$.
\end{resultbox}

%\newpage
% ==============================================================
\sectitle{Q5. Exact Line Search Formulation}{5}
% ==============================================================

\noindent\textbf{Problem.} Formulate the exact line search at $c$ along $-\nabla f(c)$
as an unconstrained optimisation in $\mu$.

\begin{resultbox}
\[
  \min_{\mu\ge 0}\;\phi(\mu) = f\!\bigl(c - \mu\,\nabla f(c)\bigr).
\]
\end{resultbox}

\reason{We are at the current point $c$ and search along the steepest-descent
        direction $-\nabla f(c)$. The scalar $\mu\ge 0$ controls how far we step.
        For a quadratic with $A\succ 0$, $\phi(\mu)$ is convex with a unique
        minimiser $\mu^*>0$, so the constraint is automatically satisfied.}

\newpage
% ==============================================================
\sectitle{Q6. Optimal Step Size $\mu^*$}{12}
% ==============================================================

\noindent\textbf{Problem.} Solve the above for $\mu^*$ and write the updated point
$d = c - \mu^*\nabla f(c)$.

Let $g = \nabla f(c) = Ac + b$ (from Q4).

\medskip
\textbf{Step 1 --- differentiate $\phi(\mu)$ via chain rule:}
\begin{align*}
  \frac{d\phi}{d\mu}
  &= \bigl(\nabla f(c-\mu g)\bigr)^\top \cdot \frac{d(c-\mu g)}{d\mu} \\
  &= \bigl(\nabla f(c-\mu g)\bigr)^\top \cdot (-g) \\
  &= -\nabla f(c-\mu g)^\top\, g = 0.
\end{align*}

\reason{Chain rule: derivative of the outer function ($\nabla f$ evaluated at $c-\mu g$)
        times derivative of the inner function ($-g$).}

\textbf{Step 2 --- evaluate $\nabla f(c-\mu g)$ using the result from Q4:}
\[
  \nabla f(c-\mu g) = A(c-\mu g)+b = (Ac+b)-\mu Ag = g - \mu Ag.
\]

\textbf{Step 3 --- set to zero:}
\begin{align*}
  -(g-\mu Ag)^\top g &= 0 \\
  -g^\top g + \mu\,g^\top Ag &= 0 \\
  \mu^* &= \frac{g^\top g}{g^\top Ag}.
\end{align*}

\reason{Expand: $-g^\top g + \mu\,g^\top Ag = 0$.
        Since $A\succ 0$, $g^\top Ag>0$, so we can divide through.}

\begin{resultbox}
\[
  \mu^* = \frac{\|\nabla f(c)\|^2}{\nabla f(c)^\top A\,\nabla f(c)},
  \qquad
  d = c - \frac{\|\nabla f(c)\|^2}{\nabla f(c)^\top A\,\nabla f(c)}\,\nabla f(c).
\]
This is the exact steepest-descent step size for a quadratic objective.
\end{resultbox}

\end{document}
